\section{The modular schemes, their fibres and their parametrizations}

\begin{statement}{B.1}{Proved}{The solution space and the parametrizations are isomorphisms}
The scheme $\mathcal X$ of A.1 is $\Ga\times\Gm$ over $\Z$.
Both maps in C3:368, (2), and C3:371, (10), are isomorphisms of
schemes over $\Z$, for the fixed modular matrices. Explicitly,
\[
 (v,a,b)\longmapsto
 (\alpha=va^{-1},\ \beta=vb^{-1},\ \mu=\det v,\ e=e(b))
\]
identifies $B'_0\times T_r\times N'_s$ with $\mathcal G^*_{r,s}$;
setting $e=1$ gives the unsigned version. They are bijective on
points over every ring, as well as isomorphisms of fppf sheaves.
There are no missed points in these parametrizations over $\Z$,
$\F_2$, $\F_3$ or $\Z/4$.
\end{statement}
\begin{proof}
Put $\eta=\left(\begin{smallmatrix}A&B\\C&D\end{smallmatrix}\right)$.
The trace condition for $s$ gives $C=B$. The trace condition for
$\zeta=\eta L_1(\mu-1)$ is $\mu D=1$. The determinant condition
then gives
\[
 \eta=\begin{pmatrix}\mu(1+B^2)&B\\B&\mu^{-1}\end{pmatrix}.
\]
Thus $\mu$ is automatically a unit, even though initially it may
have been an affine coordinate. For
$v=b_0(\mu,\mu B)$ direct multiplication gives
$\eta=[v,s]$ and $\zeta=[v,r]$. This gives a polynomial inverse
in the localized coordinate ring and proves the first assertion.

Given $(\alpha,\beta,\mu,e)$ satisfying the signed equation, use
$\eta=e[\beta,s]$ to recover the unique $v$. Set
$a=\alpha^{-1}v$ and $b=\beta^{-1}v$. Equality of commutators
implies $a\in T_r$, and $[v,s]=e[\beta,s]$ implies
$b s b^{-1}=e s$. Conversely these identities give the original
equation. All formulas are defined over the original ring; no
extension is needed. This also proves injectivity, rather than
assuming that two centralizers intersect trivially.

The output independently enumerates $\mathcal X(\F_2)$,
$\mathcal X(\F_3)$ and $\mathcal X(\Z/4)$: their cardinalities
are $2,6,8$ and the missed-point sets are empty.
The statement does not assert the analogous parametrization for
arbitrary pairs $(\rho,\sigma)$ without the modular hypotheses
discussed at C3:370--371.
\end{proof}

\begin{statement}{B.2}{Proved}{Equations, integral Groebner basis and flatness}
The final group $\mathcal M$ of C3:399, (13), and C3:479, (3),
has coordinate algebra
\[
\begin{split}
\Z[c,d,t,u,U_{11},U_{12},U_{21},U_{22},e,e',Q_2^{-1}]/
\big(&e^2-1,(e')^2-1,\\
&\det U-eQ_2,\ eQ_2-e'Q_3\big).
\end{split}
\]
It is flat of relative dimension $6$ over $\Z$. Its special
precursor $S\mathcal G^*_{r,s}$ has relative dimension $4$;
$S\mathcal M$ has relative dimension $7$. These are different
schemes and explain the corrections at scans151 and393.
\end{statement}
\begin{proof}
A.1 gives $a=cr+dI$, $b=\left(\begin{smallmatrix}u&t\\-et&eu\end{smallmatrix}\right)$.
The displayed equations are precisely the three determinant
conditions and the two sign conditions. Since $Q_2$ is a unit,
so is $Q_3$, and all three matrices are invertible. This proves
the asserted functor of points without an unmentioned open
condition.

Write $U=D(eQ_2)h$, $h\in\SL_2$. Then, as schemes,
\[
 \mathcal M\simeq\SL_2\times X,\qquad
 X=\Spec\frac{\Z[c,d,t,u,e,e',z]}
 {(Q_3-ee'Q_2,e^2-1,(e')^2-1,zQ_2-1)}.
\]
In lexicographic order
\[
 h_{11},h_{12},h_{21},h_{22},c,e,e',z,t,u,d
\]
the following five polynomials are a monic Groebner basis over $\Z$:
\[
\begin{gathered}
 h_{11}h_{22}-h_{12}h_{21}-1,\quad
 c^2+cd+d^2-ee'(t^2+u^2),\\
 e^2-1,\quad(e')^2-1,\quad z(t^2+u^2)-1.
\end{gathered}
\]
Their leading monomials
$h_{11}h_{22},c^2,e^2,(e')^2,zt^2$ are pairwise coprime.
The product criterion reduces every $S$-polynomial to zero using
only monic division, so this is an integral assertion, not a
Groebner calculation over $\Q$ misreported as one over $\Z$.
The residue classes of the standard monomials are a free
$\Z$-basis. In particular there is no integer torsion and the
scheme is flat. The same basis remains a Groebner basis in every
characteristic. The program computes it independently over
$\Q,\F_2,\F_3$ and verifies the integral monicity certificate.

For equidimensionality, $X\to\mu_2^2$ is smooth of relative
dimension $3$, as proved in B.3; therefore $\mathcal M$ has
relative dimension $6$ at every point. Finally
$S\mathcal G^*\simeq\Ga\times X$ by B.1 and the determinant
conditions, and $S\mathcal M\simeq\Ga\times\mathcal M$ as
schemes. The signed $\GL_2$ precursor $\mathcal G^*$ itself has
dimension $6$, whereas the equations with $\alpha,\beta\in\SL_2$
and no signs have dimension $4$. Scan151's marginal correction
to $4$ concerns this precursor, not the later $\mathcal M$.
\end{proof}

\begin{statement}{B.3}{Proved}{Smooth locus and the fibres at 2 and 3}
The smooth locus of $\mathcal M\to\Spec\Z$ is precisely
$\mathcal M_{\Z[1/2]}$. The closed subscheme cut out by the
Jacobian ideal for relative dimension $6$ is $V(4)$. The fibre
at $3$ is smooth and has four geometrically irreducible
components, each geometrically isomorphic to
$\SL_2\times\Gm^2\times\Ga$. The fibre at $2$ has one
geometrically irreducible component and is everywhere nonreduced;
its reduction is
\[
 \SL_{2,\F_2}\times\Ga\times\Gm\times
 R^1_{\F_4/\F_2}\Gm.
\]
Its nilradical is generated by $e-1,e'-1$, its cube is zero,
and its generic length over its reduction is $4$. The singular
locus of that fibre is the entire fibre. The unsigned subgroup
$\mathcal M^0$ is smooth over all of $\Z$.
\end{statement}
\begin{proof}
Relative to $\mu_2^2$, $X$ is the open hypersurface
$Q_3=ee'Q_2$, with $Q_2$ a unit. Its gradient in $c,d,t,u$ is
\[
 (2c+d,c+2d,-2ee't,-2ee'u).
\]
At a point of residue characteristic different from $2$, the
last two entries cannot both vanish. In characteristic $2$,
the first two become $d,c$; they cannot both vanish since
$Q_3$ is a unit. Thus this hypersurface is smooth over the sign
base in every characteristic. It is surjective after an
algebraic closure: prescribe a unit norm and solve the two
quadratic norm equations. Hence it is faithfully flat over
$\mu_2^2$. The geometric descriptions below also prove
nonemptiness explicitly.

Away from $2$ the sign base is etale. At $2$, both sign
equations have zero differential, increasing tangent dimension
by two. This proves the smooth-locus assertion. More precisely
the sign rows in the Jacobian force every maximal minor to
be divisible by $4$. The remaining gradient entries generate
the unit ideal: their ideal contains $2$, since $Q_2$ is
invertible; it then contains $c,d$ and hence the unit $Q_3$.
The $\SL_2$ gradient also generates the unit ideal, and the
localization equation has invertible $z$-derivative $Q_2$.
Consequently the minors generate exactly $(4)$.

At $3$, $Q_3=(c-d)^2$. Put $v=c-d\ne0$.
For each sign pair, the remaining equation is
$(t/v)^2+(u/v)^2=ee'$. This is a smooth norm torsor for
$\F_9/\F_3$; both values $1,-1$ occur, for example at
$(1,0)$ and $(1,1)$. The unused coordinate $d$ is affine.
This proves the description and geometric irreducibility of
each of the four components.

At $2$, put $x=e-1,y=e'-1$. Then $x^2=y^2=0$.
After reduction $Q_2=(t+u)^2$. Writing $v=t+u\ne0$ and
dividing $c,d$ by $v$ gives the norm-one equation
$a^2+ab+b^2=1$ for $\F_4/\F_2$; the remaining coordinate is
affine. This reduction is smooth and geometrically connected.
Smoothness over the Artinian sign algebra
$\F_2[x,y]/(x^2,y^2)$ proves that its nilradical is exactly
$(x,y)$, with nonzero product $xy$ and generic length $4$.
The same gradient argument with $e=e'=1$ proves smoothness
of $\mathcal M^0$ over $\Z$.

In particular the notation ``composante neutre'' for
$\ker(e,e')$ in C3:400 and419 is valid after inverting $2$,
but is not the identity component of the characteristic-$2$
fibre: that fibre is connected, and its reduction is a proper
closed subgroup. The good behaviour at $3$ is stronger than
merely smoothness away from $2$ and $3$.
\end{proof}

\begin{statement}{B.4}{Proved}{Group law and exact sequences of the modular construction}
The group law on $\mathcal M\subset T_r\times N'_s\times G$
is componentwise multiplication. With $\mathcal H$ as in A.1,
the following sequences are exact as fppf sheaves:
\[
\begin{gathered}
1\to T_s\to N'_s\xrightarrow e\mu_2\to1,\\
1\to\SL_2\to\mathcal M\to\mathcal H\to1,\qquad
1\to\Ga\to S\mathcal M\to\mathcal M\to1,\\
1\to\mathcal M^0\to\mathcal M\xrightarrow{(e,e')}\mu_2^2\to1,\\
1\to ST_r\times ST_s\times\SL_2\to\mathcal M
 \xrightarrow{(\chi,e,e')}\Gm\times\mu_2^2\to1,\\
1\to ST_r\times ST_s\to\mathcal M^0
 \xrightarrow{U}\GL_2\to1,\\
1\to ST_r\times ST_s\to\mathcal H^0\xrightarrow\chi\Gm\to1.
\end{gathered}
\]
The first three have explicit sections and are exact on points
over every ring. The other surjections need not be onto on
points over a given ring.
The larger ambient groups at C3:393 have, after the same change
of coordinates,
\[
\widetilde S\mathcal M\simeq B'_0\times T_r\times N'_s\times\GL_2,
\qquad
\widetilde{\mathcal M}\simeq\Gm\times T_r\times N'_s\times\GL_2,
\]
and a split exact sequence
$1\to L_0\to\widetilde S\mathcal M\to\widetilde{\mathcal M}\to1$.
Their relative dimensions are $10$ and $9$.
\end{statement}
\begin{proof}
Determinants and signs are characters, so their equalities are
preserved under products and inverses. The first section was
given in A.1. The section of $\mathcal M\to\mathcal H$ is
$(a,b)\mapsto(a,b,D(\det b))$; conjugation on the normal
$\SL_2$ factor is by $D(\chi)$. The section of
$S\mathcal M\to\mathcal M$ appends $D(\chi)$ in the
$B'_0$ coordinate. In our order $(v,a,b,U)$ its kernel is
$(L_0(z),I,I,I)$; the differing coordinate order in the
overwritten formulas at C3:479 must not be ignored.

All kernels follow directly from the equations. To prove
fppf surjectivity of the remaining arrows, solve prescribed
norm equations by adjoining square roots of the required
units and taking scalar elements in the respective quadratic
algebras. Such adjoining is finite locally free of rank $2$
and faithfully flat, also in characteristic $2$. Given
$(\chi,e,e')$, require $Q_2=e\chi$ and $Q_3=e'\chi$ and
take $U=D(\chi)$. This constructs the required local lift.
The same construction with $e=e'=1$ treats the final arrows.

For a pointwise failure over $\Z$, both positive quadratic
norms of integral units equal $1$, so
$(\chi,e,e')$ has image only $(1,1,1)$ and $(-1,-1,-1)$.
Over $\F_3$, $Q_3$ of a unit is always $1$, so
$\mathcal H^0(\F_3)\to\F_3^\times$ misses $-1$.
These examples disprove pointwise exactness without changing
the fppf statements. The table in B.7 gives further complete
point images for the individual commutator maps.

For the ambient groups begin with
$\GL_2\rtimes(B'_0\times T_r\times N'_s)$, where $v\in B'_0$
acts on $\GL_2$ by conjugation and the other factors act trivially.
The change $(g,v,a,b)\mapsto(v,a,b,U=gv)$ converts its product
to componentwise multiplication. The diagonal copy of $L_0$
being killed becomes $(v,I,I,I)$ with $v\in L_0$.
Its quotient replaces $B'_0$ by $B'_0/L_0=\Gm$, with splitting
$\mu\mapsto D(\mu)$. This proves the extra sequence, its pointwise
exactness and both dimensions. Imposing
$\det v=\det b=e'\det a=\det U$ gives the smaller sequence.
\end{proof}

\begin{statement}{B.5}{Proved}{The distinguished sections and the actual normalizer}
The subgroups $\mathcal M(0)$ and $\mathcal M[0]$ are
$\mathcal H$ and $\Ga\rtimes_\chi\mathcal H$ through the
embeddings in A.1. Their normalizer comparison is
\[
 \operatorname{Norm}_{\mathcal M}(L_0)
 =\{(a,b,U)\in\mathcal M:U\text{ is upper triangular}\},
\]
which properly contains $\mathcal M[0]$. The equations
$U=b$ and $U=a$ define the section $\mathcal M(s)$ and
the partial section $\mathcal M(r)$ over $e'=1$.
Conjugating the first two appropriate sections by $(I,I,r)$
defines the cycles
\[
 \{\mathcal M(0),\mathcal M(1),\mathcal M(\infty)\},\qquad
 \{\mathcal M(1/2),\mathcal M(2),\mathcal M(-1)\}
\]
with $\mathcal M(1/2)=\mathcal M(s)$ and
$\dot r(q)=1/(1-q)$, as in C3:477.
\end{statement}
\begin{proof}
Multiplying $D(\chi)$ by $L_0(z)$ gives exactly $B'_0$.
Conjugation by $D(\chi)$ sends $L_0(z)$ to $L_0(\chi z)$,
which proves the asserted semidirect product. In $\GL_2$,
the normalizer of the whole root subgroup $L_0$ is the full
upper triangular Borel: conjugating $I+zE_{12}$ and comparing
its coefficient of $z$ shows that the line spanned by
$E_{12}$ is preserved; its kernel and image are the first
coordinate line, which the conjugator must preserve.
The converse is direct multiplication. Since the $U$ projection
governs conjugation of the $\SL_2$ factor, the normalizer
formula follows. The element $(I,I,-I)$ is already in this
normalizer but not in $\mathcal M[0]$ over $\Z$.
This is the correction explicitly recognized at C3:418,
fol.510; it is not treated as a still-maintained assertion.

The two section formulas preserve multiplication and satisfy
the determinant constraint precisely on the indicated bases.
Finally $r^3=-I$ acts trivially by conjugation, so the
threefold conjugation cycles close. Their homographic labels
are the orbits of $0$ and $1/2$ under $q\mapsto1/(1-q)$.
The label $-\bar j$ names $\mathcal M(r)$ in C3:477;
the attempted additional choice $\mathcal M(-j)$ at
scans477--478 is not fixed by the source and is not silently
declared canonical here.
\end{proof}

\begin{statement}{B.6}{Proved}{The central kernels and the square-root cover}
For $\mathcal M^0$ the kernel of $U:\mathcal M^0\to\GL_2$
is $I^0=ST_r\times ST_s$, and the kernel of its composite
with $\GL_2\to\PGL_2$ is
\[
 J^0=\{(a,b,\lambda):Q_3(a)=Q_2(b)=\lambda^2\}
 \simeq ST_r\times ST_s\times\Gm.
\]
It is the centre of $\mathcal M^0$. Forgetting $\lambda$
gives an fppf central extension
\[
 1\to\mu_2\to J^0\to\mathcal H^0\to1,
\]
not the unqualified isomorphism to $\mathcal H^0$ written
in C3:400. Moreover
\[
 (J^0\times\SL_2)/\mu_2\simeq\mathcal M^0,
\]
where $z\in\mu_2$ is embedded as $((I,I,z),z^{-1}I)$.
For the full signed group the centre is
\[
 \{(a,uI,\lambda I):a\in T_r,\ u,\lambda\in\Gm,
   \lambda^2=u^2,\ (u^2/\det a)^2=1\}.
\]
\end{statement}
\begin{proof}
The first equations are immediate. Dividing $a,b$ by
$\lambda$ gives the displayed isomorphism with the product
of two norm-one groups and $\Gm$. The map $U$ is fppf
surjective by B.4, and the first two coordinates commute;
therefore an element is central in $\mathcal M^0$ exactly
when $U$ is scalar. Adjoining a square root of the common
norm proves surjectivity $J^0\to\mathcal H^0$ and identifies
its kernel. As a finite check separating the proposed
isomorphism, over $\F_3$ these groups have respectively
$6\cdot4\cdot2=48$ and $24$ points.

The multiplication map sends $((a,b,\lambda),h)$ to
$(a,b,\lambda h)$. Its kernel is the indicated diagonal
$\mu_2$, and it is locally onto by taking a square root of
$\det U$. This proves the quotient assertion as fppf sheaves.
For the full centre, projection onto $N'_s$ and $\GL_2$
is fppf surjective, by the norm-root construction. The
centre of $N'_s=T_s\rtimes\mu_2$ consists of its scalar
matrices: commuting with $T_s$ forces the sign to be $1$,
and invariance under the universal sign forces the
$s$ coefficient to vanish. This argument compares
coefficients in $R[e]/(e^2-1)$, so it is valid at $2$ as
a scheme statement. The determinant equations give exactly
the formula claimed.
\end{proof}

\begin{statement}{B.7}{Proved}{Finite point counts and complete small-ring commutator images}
For the final signed $\mathcal M$, the exact counts are
\[
\begin{array}{c|r|r|r}
R&|\SL_2(R)|&|\mathcal H(R)|&|\mathcal M(R)|\\\hline
\F_2&6&6&36\\
\F_3&24&96&2304\\
\F_5&120&384&46080\\
\F_7&336&1152&387072\\
\F_{11}&1320&5760&7603200\\
\F_{13}&2184&6912&15095808\\
\Z/4&48&192&9216\\
\Z/8&384&6144&2359296\\
\Z/9&648&2592&1679616
\end{array}
\]
For $p\ge5$,
\[
 |\mathcal M(\F_p)|=
 4(p-1)\bigl(p-(\tfrac{-1}{p})\bigr)
 \bigl(p-(\tfrac{-3}{p})\bigr)p(p^2-1).
\]
The individual maps $A\mapsto[A,r]$ and $A\mapsto[A,s]$
have the small-ring images described below; they must be
distinguished from the everywhere bijective combined
parametrizations of B.1.
\end{statement}
\begin{proof}
For a finite ring $R$ put
$n_i(v)=\#\{(x,y)\in R^2:Q_i(x,y)=v\}$ and
$E(R)=\{e\in R^\times:e^2=1\}$. Then the equations give
\[
 |\mathcal H(R)|=
 \sum_{e,e'\in E(R)}\sum_{v\in R^\times}
 n_2(v)n_3(ee'v),\qquad
 |\mathcal M(R)|=|\SL_2(R)||\mathcal H(R)|.
\]
The second equality is an equality of point sets from an
explicit section, not an assumption about exactness on
points. For $p\ge5$ each nonzero norm has $p-1$ preimages
in the split quadratic algebra and $p+1$ in the quadratic
field. This follows respectively from $xy=v$ and the
cyclic multiplicative group of the field of $p^2$ elements.
The splitting discriminants are $-1$ and $-3$, proving
the formula. The factor $|\SL_2(\F_p)|=p(p^2-1)$ follows
by choosing a nonzero first column and then one of the $p$
second columns with determinant $1$.

For $p=3$, $Q_3=(c-d)^2$, so the nonzero-value counts
are $n_3(1)=6,n_3(2)=0$ and $n_2(1)=n_2(2)=4$.
For $p=2$ they are $n_3(1)=3,n_2(1)=2$.
For $4,8,9$ the code enumerates the quadratic values and,
independently, every invertible matrix, its determinant,
its centralizer or signed-normalizer condition, and every
admissible pair. The two methods agree on each count.

For completeness, write a matrix as its four row-major
entries. The only points missed by the $\GL_2$ commutator
maps on the full numerical targets are
\[
 \begin{array}{c|cc}
 &r&s\\\hline
 \F_2&\varnothing&\{(0,1,1,0)\}\\
 \F_3&\{(2,1,2,0)\}&\varnothing\\
 \Z/4&\varnothing&\varnothing.
 \end{array}
\]
They are exactly the nonregular points. For $\SL_2$ the
missed sets are the same in the empty cases, with the
following two replacements:
\[
\begin{split}
r,\F_3:\quad\{&(0,2,1,2),(1,1,2,0),(2,0,2,2),\\
 &(2,1,0,2),(2,1,2,0)\};\\
s,\Z/4:\quad\{&(2,1,1,3),(2,3,3,3),(3,0,0,3),\\
 &(3,1,1,2),(3,2,2,3),(3,3,3,2)\}.
\end{split}
\]
The $s,\F_2$ singleton still remains. A.5 explains these
images by the regularity and norm obstructions; exhaustive
enumeration supplies every point, not only a witness.
Over $\Z$ the complete missed sets are the inequalities
in A.5. All outputs are retained.
\end{proof}

\begin{statement}{B.8}{Open}{General pairs and the fundamental realization hypothesis}
For an affine group scheme $G$ of finite presentation, fixed
sections $\rho,\sigma$, and a fixed morphism
$\ell:\Ga\to G$, the functor
\[
 (\alpha,\beta,\nu)\longmapsto
 [\alpha,\rho]=[\beta,\sigma]\ell(\nu)
\]
is represented by a closed subscheme of $G^2\times\Ga$.
Adding a specified central sign group and prescribed
characters is again a fibre-product construction. This
representability does not prove the automorphism-realization
hypothesis (37) of C3:501, fol.553, or identify every general
group of \S49 with the modular matrix model.
\end{statement}
\begin{proof}[Proved part and exact obstruction]
The two sides are morphisms to the affine separated scheme
$G$; their equalizer is a closed subscheme and represents
the functor. The equality is compatible with every base
change. In the modular $\GL_2$ case B.1 provides an explicit
inverse, and hence an actual transported group law over
every base, including $\F_2,\F_3,\Z/4$. Those finite cases
were computed above. With no signs the equal-norm group
$\mathcal H^0$ and the group $\mathcal M^0$ are already
constructed and smooth over $\Z$; adjoining the two signs
is the explicit extension of B.4.

For a topologically generated group $C=\langle\rho,\sigma\rangle$,
the assignment of the images of its generators determines
at most one continuous homomorphism. Existence requires
that every defining relation of $C$ remain true under
$\rho\mapsto\alpha\rho\alpha^{-1}$ and
$\sigma\mapsto\beta\sigma\beta^{-1}$; invertibility
additionally requires an automorphism. The single loop equation
alone supplies neither implication. B.12 gives a finite counterexample
even when the proposed images generate and the loop exponent is a
unit. This is precisely the additional assumption in C3:501 (37)(b).

For the specific completed modular group, existence has a stronger
answer. Fix
\[
 C=\widehat{\SL_2(\Z)},\qquad
 C^\pm=C\rtimes\langle\tau\rangle,\qquad
 \omega=r^3=s^2,\quad\omega^2=1,
 \tag{B.8.1}
\]
where $\tau$ is the involution induced by conjugation by
$\operatorname{diag}(-1,1)$ on the discrete matrix group and then
by continuity on its profinite completion. Thus
$\tau(s)=s^{-1}$, $\tau(r)=sr^{-1}s^{-1}$, and $\tau(\omega)=\omega$.
Let $\nu_2,\nu_3:\widehat\Z^\times\to\Z/2$ be defined by
$(-1)^{\nu_2(\mu)}=\varepsilon_2(\mu)$ modulo $4$ and
$(-1)^{\nu_3(\mu)}=\varepsilon_3(\mu)$ modulo $3$.
Define $\Sigma_C^{\rm sgn}$ to be the set of triples
$\alpha\in C^\pm$, $\beta\in C$, $\mu\in\widehat\Z^\times$
such that the parity of $\alpha$ is $\nu_3(\mu)$ and
\[
 [\alpha,r]=\omega^{\nu_2(\mu)}[\beta,s]
                 (rs^{-1})^{\mu-1}.
 \tag{B.8.2}
\]
The scalar sign is written as a power of the actual central element
$\omega$, so (B.8.2) is an equation in $C$.

Every such triple defines a unique continuous endomorphism
$u_{\alpha,\beta,\mu}:C\to C$ by
\[
 r\longmapsto \alpha r\alpha^{-1},\qquad
 s\longmapsto \omega^{\nu_2(\mu)}\beta s\beta^{-1}.
 \tag{B.8.3}
\]
Indeed, the cube of the first image and the square of the second
are both $\omega$. The images therefore satisfy the presentation
$\langle r,s\mid r^6=s^4=1,\ r^3=s^2\rangle$ proved in C.2.
They give a homomorphism from the discrete $\SL_2(\Z)$ into $C$.
The inverse image of every open normal subgroup of $C$ has finite
index, so this homomorphism is continuous for the profinite topology
and extends uniquely to its completion. This proves the assertion
without assuming automorphism realization.

Put $\ell_0=s^{-1}r$ and $\ell_1=rs^{-1}$.
Multiplying (B.8.2) gives
\[
 \begin{split}
 u_{\alpha,\beta,\mu}(\ell_0)
 &=\omega^{\nu_2(\mu)}s^{-1}
       [\beta,s]^{-1}[\alpha,r]r\\
 &=s^{-1}\ell_1^{\mu-1}r
   =\ell_0^\mu.
 \end{split}
 \tag{B.8.4}
\]
The parity of $\alpha$ and the displayed formula for $\tau(r)$
show that its image of $r$ is $C$-conjugate to
$r^{\varepsilon_3(\mu)}=r^\mu$; its image of $s$ is
$C$-conjugate to $s^{\varepsilon_2(\mu)}=s^\mu$.
This is the signed numerical condition of C3:534--535.
In the convention of B.13, $e_1=rs$ and the squared loop is
$e_1^2$. Since $rs^{-1}=\omega rs$ and $\mu-1$ is even,
\[
 (rs^{-1})^{\mu-1}=(rs)^{\mu-1}
                  =\bigl((rs)^2\bigr)^{(\mu-1)/2}.
\]
Thus the two loop conventions agree with the stated exponent;
no factor of two or central sign is discarded.

The exact group-level assertion remaining here is
\[
 \text{for every }(\alpha,\beta,\mu)\in\Sigma_C^{\rm sgn},
 \quad u_{\alpha,\beta,\mu}\text{ is surjective}.
 \tag{B.8.5}
\]
Equivalently, the two images in (B.8.3) must topologically generate
$C$. A surjective continuous endomorphism of this finitely generated
profinite group is an automorphism. To check the last implication,
there are only finitely many open normal subgroups of any given
index. Inverse image under a surjective endomorphism permutes this
finite set. Every such subgroup $N$ therefore satisfies
$u^{-k}N=N$ for some $k>0$, whence $\ker u\subset N$; intersecting
all $N$ gives trivial kernel. We have not proved (B.8.5).
B.11 proves the formal unsigned construction, and B.13 proves the
formal signed construction on the effective subset. The executed
$S_3,S_4,\SL_2(\F_3),\SL_2(\Z/4)$ cases in B.11 and B.14
verify realization in those finite cases. They do not establish
(B.8.5) for $C$. The first separate congruence-to-GT lifting
instance in our computation is the modulus-$3$ tuple of D.7;
the arithmetic comparison, its central lift, and its image and
kernel are additional questions recorded there.

There is also a precise elementary obstruction to replacing this
profinite problem by an unqualified integral fppf statement.
Let
\[
 H^{\rm const}=\varprojlim_n(\Z/n)_{\Z}
\]
be the constant profinite group scheme attached to $\widehat\Z$.
There is no group-scheme morphism
$H^{\rm const}\to\SL_{2,\Z}$ sending $1$ to $L_0(1)$.
Its coordinate ring is the filtered union of the finite constant
coordinate rings. Since $\SL_2$ is affine of finite presentation,
any morphism from $H^{\rm const}$ factors, as a scheme morphism,
through $(\Z/n)_{\Z}$ for some $n$. Its image on $\Z$-points is
then finite. A homomorphism sending $1$ to $L_0(1)$ would instead
send every integer $m$ to the distinct matrix $L_0(m)$, a
contradiction. The same argument rules out extending the standard
modular matrices from the constant profinite scheme attached to
$C$. Direct multiplication gives
$L_0(a)L_0(b)=L_0(a+b)$, which proves the asserted power identity;
D.17 records its compatible finite images.

An integral realization must consequently specify an adic or formal
construction, or a different schematic envelope, and prove its own
representability and comparison maps. Such a construction is not
supplied by the abstract profinite fibre products, and is not
identified with D.7's finite lifting question. We do not assert an
fppf quotient between undefined integral parameter functors.
\end{proof}

\begin{statement}{B.9}{Proved}{Removing a global sign is not a scheme quotient on adelic points}
Let $K$ be a profinite group with a specified continuous
homomorphism $\phi:K\to\GL_2(\widehat\Z)/\{\pm I\}$,
where the sign subgroup is diagonal and has order $2$.
Then
\[
 \widetilde K=K\times_{\GL_2(\widehat\Z)/\{\pm I\}}
                  \GL_2(\widehat\Z)
\]
is a canonical profinite central extension of $K$ by $C_2$,
equipped with a genuine matrix representation. It has a
continuous splitting exactly when $\phi$ has a continuous
homomorphic lift. On the other hand the natural map
\[
 \GL_2(\widehat\Z)/\{\pm I\}\longrightarrow
       (\GL_2/\mu_2)(\widehat\Z)
\]
has kernel $\mu_2(\widehat\Z)/\{\pm1\}$, which is infinite.
\end{statement}
\begin{proof}
The fibre product is closed in a product of profinite
groups. Projection onto $K$ is onto because the quotient
map is onto, and its kernel is the specified diagonal
$C_2$. A splitting and a lift give one another by their
two components. For the scheme quotient, exactness as
fppf sheaves says that the kernel on points is
$\mu_2(\widehat\Z)$. A unit of $\Z_p$ with square $1$
is $1$ or $-1$ since $\Z_p$ is a domain; hence
\[
 \mu_2(\widehat\Z)=\prod_p\{\pm1\}.
\]
Quotienting by the diagonal sign proves the asserted
kernel. For example the scalar equal to $-I$ at $5$ and
$I$ elsewhere is a nontrivial element of it. This gives
a concrete obstruction to restating the passage of
C3:492 as an unqualified statement about the integral
scheme quotient. Identification of the particular
$\widetilde K$ with the manuscript's
$\widetilde N_{1,1}$ needs the specified geometric
extension class; the universal pullback construction by
itself is not that identification.
\end{proof}

\begin{statement}{B.10}{Refuted}{Lead 2 with its strongest explicit modular replacement}
The final six-dimensional $\mathcal M_{r,s}$ is not the
scheme of pairs $(\alpha,\beta)$ in the preliminary
special-linear loop equation, and its $T_r,T_s$ factors
are not the norm-one tori. The precise replacements are
the schemes, dimensions and split extension in B.1--B.4.
The $\GL_2$ centralizers have dimension $2$ and the
norm-one kernels dimension $1$; the final $\mathcal M$
is an $\SL_2$ extension of the three-dimensional
equal-squared-norm group $\mathcal H$.
\end{statement}
\begin{proof}
The coordinate algebras and inverse maps prove each
identification. A dimension-$4$ precursor cannot be
isomorphic to the flat dimension-$6$ final group.
Over an algebraic closure in characteristic different
from $2,3$, $T_r,T_s$ are $\Gm^2$, whereas their
norm-one kernels are $\Gm$; dimensions alone disprove
their identification. Their degeneration at the
ramified primes is explicitly computed in B.3 and in
the norm equations, so no torus terminology is extended
across those primes without qualification.
\end{proof}

\begin{statement}{B.11}{Proved}{The general unsigned construction, with its hypotheses}
Let $C$ be a discrete group, or a topologically finitely generated
profinite group, generated by $\rho,\sigma$. Put
$L_0=\overline{\langle\sigma^{-1}\rho\rangle}$ (omit the closure
in the discrete case), and let $B\subset\Aut(C)$ be the full
subgroup preserving both conjugacy classes and normalizing $L_0$. Write
\[
 K=L_0\cap Z(C),\quad \gamma=B/\operatorname{Int}(L_0),\quad
 G=(C\rtimes B)/\{(\ell^{-1},\operatorname{Int}\ell):\ell\in L_0\}.
\]
All homomorphisms below are continuous in the profinite case.
The map $B\to G$ is injective, the kernel of $C\to G$ is exactly
$K$, and
\[
 1\longrightarrow C/K\longrightarrow G\longrightarrow\gamma
 \longrightarrow1
\]
is exact. Assume now $K=1$. Put $Z_\rho=C_G(\rho)$,
$Z_\sigma=C_G(\sigma)$, $C_\rho=C_C(\rho)$ and
$C_\sigma=C_C(\sigma)$. The arrows $Z_\rho,Z_\sigma\to\gamma$
are onto, and the groups
\[
 T^0=Z_\rho\times_\gamma Z_\sigma,\quad
 M^0=T^0\times_\gamma G,\quad SM^0=B\times_\gamma M^0
\]
satisfy
\[
\begin{gathered}
1\to C_\rho\times C_\sigma\to T^0\to\gamma\to1,\\
1\to C_\rho\times C_\sigma\times C\to M^0\to\gamma\to1,\\
1\to L_0\to SM^0\to M^0\to1,\\
1\to L_0\times C_\rho\times C_\sigma\times C
 \to SM^0\to\gamma\to1.
\end{gathered}
\]
These are the corrected fibre-product sequences of C3:505--513,
\S49.II, fols.555--559 (parallel overwritten folios504--512).
If the realization map $\Phi:B\to\Sigma$ of C3:503, (9)--(10),
is bijective, the parametrizations (34), (38) are bijective,
with the inverse and transported group law described below.
\end{statement}
\begin{proof}
We use $(c,u)(c',u')=(c\,u(c'),uu')$ in the semidirect product.
The map $i(\ell)=(\ell^{-1},\operatorname{Int}\ell)$ is an
injective homomorphism. Its image centralizes $C$, and conjugation
by $u\in B$ sends $i(\ell)$ to $i(u\ell)$, proving normality.
In the profinite case it is closed because $L_0$ is compact.
The automorphisms of a finitely generated profinite group form a
profinite group: there are only finitely many subgroups of each
index, their intersections give a cofinal family of characteristic
open subgroups, and compatible automorphisms of their finite
quotients are exactly continuous automorphisms of $C$.
The conditions on $B$ are closed, since conjugacy classes are
compact and $u(L_0)=L_0$ is a closed condition (use both $u$ and
$u^{-1}$). Thus all quotients have the asserted topology.

An element $(c,1)$ belongs to $i(L_0)$ exactly when
$c=\ell^{-1}$ and $\operatorname{Int}\ell=1$, proving the kernel
formula. An element $(1,u)$ can belong to $i(L_0)$ only for
$\ell=1,u=1$, proving injectivity of $B$. Projection onto
$B/\operatorname{Int}(L_0)$ has kernel precisely the image of $C$.
The construction is universal for two homomorphisms from $C,B$
respecting the given action and identifying $\ell\in L_0$ with
$\operatorname{Int}\ell\in B$: this follows from the universal
property of a semidirect product and its displayed kernel.
For example if $C=C_6$ with generators of orders $3,2$, then
$B=1,L_0=C,K=C,G=1$. Omitting $K$ would be false; the
manuscript already corrects the injectivity arrow at scan505.

Given $u\in B$, choose $a\in C$ with
$u(\rho)=a\rho a^{-1}$. Then $a^{-1}u\in Z_\rho$ maps to the
same element of $\gamma$ as $u$. This proves surjectivity; the
proof for $\sigma$ is identical. Their kernels are exactly
$C_\rho,C_\sigma$. Taking fibre products proves every displayed
kernel and surjectivity. The projection $SM^0\to M^0$ is the
pullback of $B\to\gamma$, whose kernel is $L_0$ because $K=1$.
These arguments prove the corresponding statements for group
objects in a sheaf topos when the displayed maps are epimorphisms
and the quotients exist there. They do not assert representability
of every general quotient.

Here $\Sigma$ consists of $(\zeta,\eta)\in C^2$ such that
$\zeta\rho$ is conjugate to $\rho$, $\eta\sigma$ is conjugate
to $\sigma$, and $\eta^{-1}\zeta\in L_1=\rho L_0\rho^{-1}$.
Assume that $\Phi(u)=(u(\rho)\rho^{-1},u(\sigma)\sigma^{-1})$
is bijective. From a loop pair $(\alpha,\beta)\in G^2$, recover
\[
 u=\Phi^{-1}([\alpha,\rho],[\beta,\sigma]),\qquad
 a=\alpha^{-1}u\in Z_\rho,\quad b=\beta^{-1}u\in Z_\sigma.
\]
Conversely $(u,a,b)\mapsto(ua^{-1},ub^{-1})$ gives that pair.
This proves (34). Its restriction to $C^2$ is exactly the
fibre product $B\times_\gamma Z_\rho\times_\gamma Z_\sigma$,
proving (38). Transporting componentwise multiplication gives
\[
 (\alpha',\beta')(\alpha,\beta)
   =(u'(\alpha)\alpha',u'(\beta)\beta'),
\]
where $u'$ acts by conjugation in $G$ and by its given
automorphism on $C$. This proves associativity without assuming
that the loop-pair set was already a group.

The symbols $ua^{-1},ub^{-1}$ at scan510, (38), must therefore
be $\alpha,\beta$, not the inserted $\zeta,\eta$, which are
their commutators. In scan511, (44), the second $SZ_\rho$ must
be $SZ_\sigma$. In scan509, (31), the final exponent on $\rho$
must be $-1$. These are mathematical readings, not claims of
final palaeographic settlement.

The finite program constructs all automorphisms by checking
generator images against every Cayley-graph edge. For
$(C,\rho,\sigma)=(S_3,(123),(12))$ and
$(S_4,(123),(34))$ it proves $\Sigma=\Phi(B)$ and obtains
\[
\begin{array}{c|rrrrrrrr}
C&|B|&|L_0|&|\gamma|&|G|&|Z_\rho|&|Z_\sigma|&|M^0|&|SM^0|\\\hline
S_3&2&2&1&6&3&2&36&72\\
S_4&8&4&2&48&6&8&576&2304
\end{array}
\]
The algorithms and exact assertions are in
\texttt{AB\_general\_finite.py}, with its JSON and text output.
\end{proof}

\begin{statement}{B.12}{Refuted}{The loop equation does not imply realization}
In $C=\mathrm{PSL}_2(\F_{13})$ there are two generating pairs
$(\rho,\sigma)$ and $(\rho',\sigma')$ whose corresponding
entries are conjugate, of respective
orders $3,2$ with
\[
 (\sigma')^{-1}\rho'=(\sigma^{-1}\rho)^2,\qquad
 |\sigma^{-1}\rho|=7,
\]
yet no homomorphism $C\to C$ sends the first pair to the second.
Thus the realization assumption of C3:503, (10)(b), cannot be
deleted, even when the loop multiplier is a unit. The general
replacement is B.11 with realization verified separately, as
it is for the matrix parametrization in B.1.
\end{statement}
\begin{proof}
All matrices denote projective classes over $\F_{13}$. Take
\[
\begin{gathered}
\rho=\begin{pmatrix}1&2\\5&11\end{pmatrix},\quad
\sigma=\begin{pmatrix}0&1\\12&0\end{pmatrix},\quad
\ell=\sigma^{-1}\rho=\begin{pmatrix}5&11\\12&11\end{pmatrix},\\
\rho'=\begin{pmatrix}2&5\\9&10\end{pmatrix},\quad
\sigma'=\begin{pmatrix}1&4\\6&12\end{pmatrix}.
\end{gathered}
\]
Conjugators are
$\alpha=\left(\begin{smallmatrix}6&6\\10&8\end{smallmatrix}\right)$
and $\beta=\left(\begin{smallmatrix}6&8\\7&3\end{smallmatrix}\right)$.
Direct multiplication verifies their determinants, conjugacies,
the loop identity and the orders. Breadth-first enumeration
from each pair produces all $13(13^2-1)/2=1092$ projective
matrices, proving generation. But
\[
 [\sigma,\rho]=\begin{pmatrix}3&12\\12&5\end{pmatrix}
 \quad\hbox{has order }7,\qquad
 [\sigma',\rho']=\begin{pmatrix}2&0\\11&7\end{pmatrix}
 \quad\hbox{has order }6.
\]
The relation $[\sigma,\rho]^7=1$ is not preserved, so no
homomorphism with these images exists.

With $\zeta=\rho'\rho^{-1}$ and $\eta=\sigma'\sigma^{-1}$,
one also checks $\eta^{-1}\zeta=\rho\sigma^{-1}=\varepsilon_1$,
exactly scan503's condition. The exponent $2\bmod7$ is a unit.
For the profinite formulation it lifts to $\widehat\Z^\times$
by taking its $7$-adic component $2$ and all other components
$1$; the integer $2$ is not itself a profinite unit. Every finite
equality, order and generation assertion above was executed by
the attached program. Its output retains an independent failed
relation from the Cayley-graph consistency algorithm as well.
\end{proof}

